% !TEX root = ../main.tex
% =========================================================
% RÉSUMÉ (FR) / ABSTRACT (EN)
% Même contenu dans les deux langues. Cible : ~150-200 mots par version.
% =========================================================
\thispagestyle{empty}

\begin{center}
    {\Large \bfseries Résumé}
\end{center}
\vspace{0.4cm}

\justifying
Ce projet de fin d'études porte sur la conception et le développement d'une
plateforme web de gestion des mobilités internationales étudiantes, réalisé au
sein du service des Relations Internationales de Toulouse INP-ENSEEIHT. La
gestion de ces mobilités repose aujourd'hui sur des données dispersées entre
plusieurs systèmes (MoveON, Pégase, Eudonet) et sur de nombreuses tâches
manuelles, sources de perte de temps et de manque de visibilité. La plateforme
proposée centralise ces informations et couvre les mobilités académiques
sortantes, les stages internationaux, les mobilités complémentaires et l'accueil
d'étudiants étrangers. Elle automatise en particulier la gestion des
candidatures et l'affectation des étudiants aux destinations, à l'aide d'un
algorithme d'affectation et de machines à états sécurisant le cycle de vie des
entités. Le projet a été mené selon la méthodologie Scrum sur six sprints, avec
une architecture fondée sur Django Ninja et Next.js, une chaîne d'intégration
continue, et un moteur de recommandation de destinations reposant sur un modèle
d'apprentissage supervisé.

\vspace{0.4cm}
\noindent\textbf{Mots-clés :} mobilité internationale, plateforme web, Django,
Next.js, algorithme d'affectation, machine à états, système de recommandation,
apprentissage automatique, Scrum.

\vspace{1.2cm}

\begin{center}
    {\Large \bfseries Abstract}
\end{center}
\vspace{0.4cm}

This graduation project deals with the design and development of a web platform
for managing international student mobility, carried out within the
International Relations office of Toulouse INP-ENSEEIHT. Today, the management of
these mobilities relies on data scattered across several systems (MoveON,
Pégase, Eudonet) and on many manual tasks, which leads to wasted time and poor
visibility. The proposed platform centralises this information and covers
outgoing academic mobility, international internships, complementary mobility and
the hosting of foreign students. In particular, it automates application
management and the assignment of students to destinations, using a matching
algorithm and state machines that secure the life cycle of the entities. The
project was conducted following the Scrum methodology over six sprints, with an
architecture based on Django Ninja and Next.js, a continuous integration
pipeline, and a destination recommendation engine based on a supervised machine
learning model.

\vspace{0.4cm}
\noindent\textbf{Keywords:} international mobility, web platform, Django, Next.js,
matching algorithm, state machine, recommender system, machine learning, Scrum.
